% !PRINTED 95
\weberpage{95}
%<<CONT
to be set, and accordingly [by (6), (7), (10)] we obtain
\[
\begin{gathered}
\omega'_1 = \frac{a+b\omega}{n}, \quad \omega'_2 = \frac{c+\partial\omega}{n}, \\
a'_1 = m'b, \quad a'_2 = m'\partial.
\end{gathered}
\]
The function $T'(u, \omega'_1, \omega'_2)$ therefore satisfies the conditions (1):
\[
\begin{gathered}
T(u+\omega'_1) = (-1)^{g'_1} e^{-\pi i m'b(2u+\omega'_1)} T(u), \\
T(u+\omega'_2) = (-1)^{g'_2} e^{-\pi i m'\partial(2u+\omega'_2)} T(u),
\end{gathered}
\]
% [sic] German source prints unprimed T on both rows although the defined function is T'.
and from this it follows that the product
\[
e^{\frac{\pi i m'bu^2}{\omega'_1}} T'(u, \omega'_1, \omega'_2)
\]
is a $\Theta$-function of the order $m'$, with the arguments
\[
\frac{u}{\omega'_1}, \quad \frac{\omega'_2}{\omega'_1}
\]
and the characteristic $(g'_1, g'_2)$. We can express this in the
equation:
\weq{12}{e^{-\frac{\pi i m' n b u^2}{a+b\omega}} \Theta^{(m')}_{g'_1,g'_2}\left(\frac{nu}{a+b\omega},\ \frac{c+\partial\omega}{a+b\omega}\right) = \Theta^{(m'n)}_{g_1,g_2}(u,\omega),}
in which the characteristics are determined by (11). The means for the
representation of these functions are contained in § 21.

We denote the transformation of $T$ and $T'$ by a single
letter $S$, or by $(T', T)$, thus:
\[
S = (T', T).
\]

If $S'$ denotes a second transformation by which $T'$ passes into $T''$,
that is,
\[
S' = (T'', T'),
\]
then we can derive from these a new transformation $S''$ by which
$T$ passes into $T''$. This one is said to be composed from $S$ and $S'$,
and is denoted thus:
\[
S'' = S'S
\]
or
\[
(T'', T) = (T'', T')(T', T).
\]
In this composition the commutative law does not in general
hold; thus $SS'$ is different from $S'S$. The associative law,
however, does hold, and it is expressed in the formula:
\[
(T''', T'')\big[(T'', T')(T', T)\big] = \big[(T''', T'')(T'', T')\big](T', T) = (T''', T).
\]
